
\documentclass[12pt]{article}

\usepackage{amsmath, amssymb, amsthm}
\usepackage{enumitem}
\usepackage{hyperref}
\usepackage{geometry}
\geometry{a4paper, margin=1in}

% Define theorem styles
\newtheorem{theorem}{Theorem}
\newtheorem{lemma}[theorem]{Lemma}
\newtheorem{proposition}[theorem]{Proposition}
\newtheorem{corollary}[theorem]{Corollary}
\newtheorem{definition}[theorem]{Definition}

% Document title and author
\title{Petri nets: Liveness is monotonic in neutral cycles}
\author{Heorhii Lopattn}
\date{\today}

\begin{document}
\newcommand{\fire}[1]{\xrightarrow{#1}}
\newcommand{\fireb}[1]{\xleftarrow{#1}}

\newcommand{\firee}[1]{\,[#1\rangle}

\maketitle

\section*{Solution}
Assume that the marking $M$ is live, and $M \le M'$. 
We will prove the following statement that directly implies that $M'$ is live: \\
For any $\sigma$, there exist $\sigma_1, \sigma_2$ s.t. $M' \fire \sigma M'' \fire {\sigma_1} M_1' \ge M_1 \fireb {\sigma_2} M$. \\
We will use induction on $\sigma$. \\
The base case is obvious, as $M \le M'$.
In the induction step we let $\sigma = \tilde \sigma\, t_i$ where $t_i$ is a transition from $p_i$ to $p_{i+1 \text{ mod } n}$. 
From the inductive assumption we have $\tilde {\sigma_1}, \tilde {\sigma_2}$, s.t. $M' \firee {\tilde \sigma} \firee {\tilde \sigma_1} \ge M \firee {\tilde \sigma_2}$ \\
There are two diferent cases: \\
If $t_i$ appears in $\tilde {\sigma_1}$, it is interchangable with the rest of the transition, so we can rewrite  $\tilde {\sigma_1} = t_i \, \tilde {\sigma_1}'$. This means we can take $\sigma_1 = \tilde {\sigma_1}'$ and $\sigma_2 = \tilde {\sigma_2}$.\\
If $t_i$ doesn't appear in $\tilde {\sigma_1}$, observe that $M'\firee{\tilde{\sigma_1}}\firee {t_i} = M'\firee {t_i}\firee{\tilde{\sigma_1}}$. \\
Since by inductive hypothesis, $M'\firee {\tilde \sigma } \firee {\tilde {\sigma_1}} \tilde {M_1} \ge M \firee {\tilde {\sigma_2}} \tilde {M}'_1$, we use liveness of $\tilde M'_1$ to obtain a firing sequence $\alpha = \alpha_1\, t_i\, \alpha_2$. Obviously it also applies to $\tilde {M_1}$, as it is greater. It holds that 
$$\tilde {M_1} \firee \alpha \ge \tilde M'_1 \firee \alpha$$
$$\iff M'\firee {\tilde \sigma } \firee {\tilde {\sigma_1}} \firee \alpha \ge M \firee {\tilde {\sigma_2}} \firee \alpha$$
$$\iff M'\firee {\tilde \sigma t_i} \firee {\tilde {\sigma_1}} \firee {\alpha_1 \, \alpha_2} \ge M \firee {\tilde {\sigma_2}} \firee \alpha$$

And we take $\sigma_1 = \tilde {\sigma_1} \, \alpha_1 \, \alpha_2$ and $\sigma_2 = \tilde {\sigma_2}\, \alpha$


\end{document}
